% Storyline 3, paragraph 3: from information to judgement.
Existing remedies pursue both routes mainly by acting on the information that observers hold. They make it shared, through public institutions, gossip or the synchronization of private opinions\cite{radzvilavicius2021,kawakatsu2024,murase2024synchronization,tham2026}, or richer, through scores, records or the aggregation of several observations\cite{schmid2023,clark2020,michelmata2024}. Better information reduces erroneous and disputed judgements but does not eliminate them, because observation and assessment remain noisy\cite{hilbe2018,fujimoto2023}. What remains is a problem of judgement, namely how a standing verdict about another person should be revised as new evidence arrives. Any revision rule faces a trade-off between condemning cooperators on disputed evidence and excusing defectors whose misconduct is not yet conclusive. Human institutions often manage this trade-off through procedure. They escalate sanctions only with repeated violations\cite{ostrom1990,cox2010} and restore an offender's standing only after a period of good conduct\cite{blumstein2009}, and impressions of other people are likewise revised asymmetrically\cite{skowronski1989,slovic1993,mendesiedlecki2013diagnostic,siegel2018}. Models of indirect reciprocity, by contrast, mostly revise reputations through a single symmetric rule, whether the latest assessment, a score or a fixed window of observations\cite{ohtsuki2004,hilbe2018,schmid2023,michelmata2024}. How the rule by which private judgements are revised shapes the punishment dilemma, and how such a rule should be designed, remains open.
